\documentclass[11pt,a4paper]{article}
\usepackage[T1]{fontenc}
\usepackage{lmodern,amsmath,amssymb,amsthm,mathtools,booktabs,microtype}
\usepackage[margin=25mm]{geometry}
\usepackage[colorlinks=true,linkcolor=blue,citecolor=blue,urlcolor=blue]{hyperref}
\newtheorem{theorem}{Theorem}[section]
\newtheorem{lemma}[theorem]{Lemma}
\newtheorem{corollary}[theorem]{Corollary}
\newcommand{\Q}{\mathbb Q}\newcommand{\Z}{\mathbb Z}
\setlength{\emergencystretch}{2em}
\title{Local Precision and Height Contradictions\ in Elliptic Integer Definitions}
\author{}\date{13 November 2026}
\begin{document}\maketitle
\begin{abstract}
We isolate strict inequalities connecting contact-height exponents,
local comparison precision, and a Lucas growth argument. For the
height exponent $24$ in Sun's twenty-eight-unknown construction,
orders $49,13,25$ suffice in place of $64,16,48$, while retaining the
original degree-$64$ logarithm truncation. The conclusion keeps the
original arithmetic hypotheses and the original number of free variables.
We also record a general valuation lemma for iterated quadratic norm
expressions and sharpen the final exponents of the proposed
214-coordinate quartic, conditional on its stated height supplier.
No new numerical defining threshold or variable record is asserted.
\end{abstract}

\section{A uniform height criterion}
Here $h$ is the logarithmic rational height. All constants in the
following theorem are independent of the fixed multiplier $M$,
the input, and the witnesses. This uniformity is an explicit hypothesis.
\begin{theorem}\label{thm:height}
Let $q>0$ and let $k_D,k_I,k_S$ be positive integers satisfying
\[
                    k_D>2q,\qquad k_I>q/2,\qquad k_S>q.
\]
Suppose a family of rational certificates supplies $\beta\in\Q$,
integers $n,s,t,z$, real $B\ge1$, and $c=|s-1|\ge M$, with $B\ge c-1$.
Suppose its integer residual $\Delta$ and, when $\beta\ne n$,
$J=|\operatorname{num}(\beta-n)|\ge1$ satisfy
\begin{align*}
 |\Delta|&\le C_D B^4,\\
 \Delta\ne0&\Longrightarrow B^2\le C_0(1+|\Delta|^{q/k_D}),\\
 \beta\ne n&\Longrightarrow B^2\le C_0(1+J^{q/k_I}),
       \qquad J\le2B e^{h(\beta)},\\
 h(\beta)&\le C_S c^{q/k_S},\\
 \Delta=0&\Longrightarrow\log B>c-4.
\end{align*}
Then one sufficiently large fixed $M$ forces $\Delta=0$ and $\beta=n$
for every certificate in the family.
\end{theorem}
\begin{proof}
For $\Delta\ne0$, the first two inequalities imply
$B^2\le C_0(1+C_D^{q/k_D}B^{4q/k_D})$.
Since $4q/k_D<2$, this bounds $B$ by a constant independent of $M$.
Choosing $M$ larger than this bound plus one excludes that case.

Put $u=q/k_I<2$. If $\beta\ne n$, then $J\ge1$ gives
$B^2\le2C_0J^u\le2^{u+1}C_0B^u e^{u h(\beta)}$.
Therefore
\[
 \log B\le C_1+\frac{q}{2k_I-q}h(\beta)
 \le C_1+\frac{qC_S}{2k_I-q}c^{q/k_S},
\]
where $C_1$ is fixed. The exponent $q/k_S<1$ makes the last expression
smaller than $c-4$ for all sufficiently large $c$. This contradicts
the last hypothesis, since $c\ge M$.
\end{proof}
No coprimality between different comparison radicals is used in this
criterion. Their individual divisibilities are sufficient to obtain
the displayed height estimates.

\section{Valuations in quadratic norm expressions}
\begin{lemma}\label{lem:tree}
Let $v$ be a discrete valuation with value group $\Z$, and let
$v(r)=R>0$ be odd. Evaluate a finite binary tree whose internal operation
is $N_r(u,w)=u^2+rw^2$. If its root has valuation greater than $-R/2$,
assign $c=1/2$ at the root and assign $c/2,(c+1)/2$ to the left and
right children of a node assigned $c$. Every leaf assigned $c$ has
valuation greater than $-cR$, and $0<c<1$ throughout the finite tree.
In particular a leaf $G/r^k$, with integer $k\ge1$, satisfies $v(G)\ge k$.
\end{lemma}
\begin{proof}
When both summands are nonzero, their valuations $2v(u)$ and
$R+2v(w)$ have opposite parity. Thus the valuation of their sum is
the minimum, not a cancellation value. Zero summands have valuation
$+\infty$ and cause no difficulty. A lower bound $-cR$ at a node
therefore gives $v(u)>-cR/2$ and $v(w)>-(c+1)R/2$.
Induction proves the assertion. Finally
$v(G)>(k-c)R>k-1$, since $R\ge1$ and $c<1$.
\end{proof}
For the seven-leaf tree in \cite[Section 5]{Sun}, the coefficients are
\[
       1/16,\ 9/16,\ 5/16,\ 13/16,\ 3/8,\ 7/16,\ 15/16.
\]
The lemma is not restricted to that tree or to the orders $1,48$.
If the constant leaf $1$ lies on the all-left path and every other leaf
is integral and zero modulo the maximal ideal, the root is one modulo
that ideal. This follows by induction when $r$ is integral and supplies
the corresponding nonzero condition in completeness arguments.

\section{The precision change in Sun's construction}
Retain all arithmetic statements consumed by \cite[Sections 2--6]{Sun},
including the group description, the exponent-$24$ contact-height bound,
the pole obstruction and its prepared-input converse, the simultaneous
Gaussian prime-pattern theorem, and the stated real and finite-place
preparation statements. These are premises of the application here,
not assertions established by our finite checker.

\begin{corollary}\label{cor:Sun}
In Sun's construction, keep the logarithm truncations $\ell_{64}$ and
$\ell_{64}^{\sigma}$, all nine point roles, the three ordinary pole
predicates, and the scalar pole predicate. Change the orders of the
four ratio differences and the Lucas residual from $64$ to $49$;
change the input comparison order from $16$ to $13$; and change the two
scalar ratio-component orders from $48$ to $25$.
Recompute the normalizing powers with those orders, and choose the
fixed multiplier using the resulting height constants.
Under the same arithmetic premises, the resulting source defines exactly
$\Z$ and has the same $28$ existential unknowns, or $29$ total coordinates.
\end{corollary}
\begin{proof}
The reciprocal-abscissa tests keep order one. They place the appropriate
points in the two formal groups. At least one anchor index is a unit
at a consumed prime. The corresponding exact ratio and truncated ratio
are integral. A positive-order ratio-difference test makes the other
truncated ratio integral; the formal-logarithm unit then makes the
other exact ratio integral. Only after this step is relative precision
converted to absolute precision. The unchanged truncation still has
precision $64$, so subtracting the two anchor equations proves
$L_i\equiv n_i\pmod {p^{49}}$ from the revised ratio-difference tests.
The Lucas and input tests give $p^{49}\mid\Delta$ and $p^{13}\mid J$.

The ordinary norm--trace implication is independent of the chosen
positive integer orders. More explicitly, at an odd contact of depth
$\nu>0$, the factors have valuations
\[
 \min\{0,\,2v(D_i)+(1-2k_i)\nu\}.
\]
They cannot cancel. The trace and pole obstruction give a place at
which every such summand is greater than $-\nu/2$. Hence
$v(D_i)>(k_i-3/4)\nu\ge k_i-3/4$, and $v(D_i)\ge k_i$.
This proves the revised individual radical divisibilities for all three
scales, without any coprimality claim between them.

Lemma~\ref{lem:tree} gives order $25$ for the scalar ratio components.
The formal logarithm preserves the valuation of the scalar ratio,
so its index identity yields $M_{\rm sc}^{25}\mid(s-1)$.
The unchanged height supplier gives
\[
 B^2\ll1+|\Delta|^{24/49},\qquad
 B^2\ll1+J^{24/13},\qquad h(\beta)\ll c^{24/25}.
\]
The implied constants remain fixed before $M$ is chosen. Theorem~\ref{thm:height}
now applies, since $49>48$, $13>12$, and $25>24$. In particular the
input estimate is $\log B\le C+12h(\beta)$.
The original Lucas implication and $\beta=2X^2+3$ give soundness:
integrality of $\beta$ forces the reduced denominator of $X$ to have
square dividing $2$, and hence to be one.

For completeness, choose the root indices before the contact primes,
exactly as in the original construction. The logarithm ratios at the
selected places still have precision $64$, which is more than every
revised requirement. Thus the nonconstant terms in the ordinary norm
factors have positive valuation. At a fixed preparation prime $q$, with
$\nu=v_q(w^2)$ and $\alpha_j=v_q(A_j)$, the lower bound remains
$3\nu+(10k_i-8)\alpha_j$, which tends to infinity with $\nu$ for each
fixed scale and order. The same square-denominator and trace-unit
arguments apply.

For the scalar tree, replace the local requirement
$v_q(s-1)-v_q(h)>48v_q(r_0)+1$ by the analogous requirement with $25$.
The coefficientwise preparation at primes above $2$ is retained.
The real bounds are also uniform: if $R=|r_0|\ge1$ and the scalar
coefficients are bounded by $A$, the bottom pair involving the ratio
now has absolute value at most $2A^2/R^{49}$. Its parent is bounded
by $(A^2+4A^4)/R^2$, just as before. Thus the same form of input-independent
real bound supplies the prepared pole inputs. All subsequent choices
are made in the original order.

Only integer exponents in rational formulas have changed. The finite
point and pole-root choices and the twenty-seven free rational
parameters are unchanged. The same all-specializations elimination
argument adds one existential coordinate. No generic-open restriction
or extra free coordinate has been introduced.
\end{proof}
The three orders are the least positive integers satisfying the three
strict inequalities of Theorem~\ref{thm:height} for $q=24$.
They are not lower bounds for every conceivable argument. We retain
$\ell_{64}$ deliberately: reducing its degree is not needed here.
No numerical total-degree bound or expanded coefficient table for the
modified polynomial is claimed.

\section{Keeping the stronger exponent in the 214-coordinate argument}
Consider the construction at upstream revision
\texttt{c84bbfa4dbe8f32fd3bf29df029f7b7849d36202}, specifically
\nolinkurl{research/quartic214/proofs/uniform-existence.md}, P08--P12.
This section is a consequence of its explicitly stated input estimates,
not a new verification of their geometric or analytic foundations.

\begin{lemma}\label{lem:blocks}
Suppose two contact blocks have radicals $M_1,M_2\ge1$, disjoint outside
a fixed finite set, and both give
$h(t)-C_0\le C_1M_i^\alpha$ when $h(t)>C_0$.
Then $h(t)\le C_0+C_1(M_1M_2)^{\alpha/2}$.
Consequently, if two comparison scales separately have radical
$M_j$ with $M_j^k\mid J\ne0$, and the aggregate trace comparison is
$B^2\le A h(T)+A'$, then
\[
 B^2\le C|J|^{\alpha/(2k)}+A'
\]
for a fixed $C$, after adjusting the fixed additive constants.
\end{lemma}
\begin{proof}
Multiply the two positive inequalities and take square roots.
For the scales $t_0=\lambda$, $t_1=T\lambda$, use
$h(T)\le h(t_0)+h(t_1)$ and the two separate divisibilities.
The case $h(t)\le C_0$ is included by the additive constant.
\end{proof}
The quoted supplier gives $\alpha=2103/1024$; replacing it by $7/3$
is unnecessary. Lemma~\ref{lem:blocks} gives
\[
 h(t)\ll M^{2103/2048},\qquad
 \Delta\ne0\Longrightarrow B^2\le C'B^{701/512}+C'',
\]
because $k=3$ and $|\Delta|\le14B^4$.
The power needed when solving the last inequality is
\[
 \frac1{2-701/512}=\frac{512}{323},
\]
in place of $9/4$. For the input comparison, $k=1$ and
$J\le2B e^{h(\beta)}$ give, above the fixed additive-height threshold,
\[
 \log B\le C+\frac{2103}{1993}h(\beta),
\]
in place of $C+(7/5)h(\beta)$.
These improvements change no coordinate or quadratic row. They do not
remove the ineffective aggregate-height constant. Since $\alpha>2$,
order two still does not satisfy the strict Lucas inequality in this
argument.

\section{A multiplicity-preserving height average}
The product construction in the proof of \cite[Lemma 16.2(1)]{GR}
keeps repeated simple factors. The following elementary consequence
explains the precise use needed by the quartic height argument.
\begin{lemma}
Let $J$ be a $g$-dimensional abelian variety over a number field $L$,
and let $A$ be simple over $L$, of dimension $a$.
Suppose $J$ contains an $L$-subvariety $L$-isogenous to $A^r$, $r\ge1$.
Assume there is an $L$-isogenous variety $J_0$, isomorphic over $L$ to
a product of simple varieties, with $h_F(J_0)\le h_F(J)$.
Assume also $h_F(B)\ge-c\dim B$ for all remaining product factors,
with the same height normalization. Then some $L$-simple factor
$C$, $L$-isogenous to $A$, satisfies
\[
                h_F(C)\le\frac{h_F(J)+cg}{r}-ca.
\]
\end{lemma}
\begin{proof}
Semisimplicity up to $L$-isogeny gives at least $r$ occurrences of the
simple isogeny class of $A$ in $J_0$. Select $r$ occurrences, even if
some factors are isomorphic, and write $J_0=C_1\times\cdots\times C_r\times E$.
Then $\dim E=g-ra$. Height additivity and the lower bound give
$\sum_i h_F(C_i)\le h_F(J)+c(g-ra)$. At least one factor is no larger
than the average. No invariance under arbitrary isogeny is used.
\end{proof}
For surfaces and $c=3/2$ this retains a subtractive term $3$ compared
with the coarser bound $(h_F(J)+3g/2)/r$. One must still transfer height
from the selected factor to the original $A$ by an isogeny in the
correct direction and pay its half-logarithmic degree cost.

\section{Verification and scope}
The exact checker verifies the rational exponents, the norm-tree
coefficients, and the valuation residue classes of the revised quartic
factors. It rejects the three boundary precision choices $48,12,24$
for the strict criterion. Its small division-polynomial and scalar
congruence tests are explicitly finite evidence, not universal proofs.
The source constants in the two applications are different; they are
not identified. The product-height result in the companion paper
and the sum-height input used by the quartic are likewise distinct.

\begin{thebibliography}{9}
\bibitem{Sun} Zhi-Wei Sun, \emph{$\Z$ is Diophantine over $\Q$ with
28 unknowns}, arXiv:2611.14605v1, 10 November 2026,
\url{https://arxiv.org/html/2611.14605v1}.
\bibitem{GR} \'Eric Gaudron and Ga\"el R\'emond,
\emph{Nouveaux th\'eor\`emes d'isog\'enie}, M\'emoires de la Soci\'et\'e
Math\'ematique de France 176 (2023), DOI: 10.24033/msmf.484.
The directly inspected text for the product construction is the authors'
September 2022 version, Lemma 16.2(1), printed p.~104,
\url{https://lmbp.uca.fr/~gaudron/art20.pdf}.
The final published page comparison is not certified here.
\end{thebibliography}
\end{document}
